%% kapa: exam=1 round=27 year=2016 period=2 form=A subject=accounting subject_ko=회계학 seq=25 number=65
%% source: Q-Net 감정평가사 기출문제 (cst00309) artlSeq=5169402 file=extracted/1008/5169402/제27회 감정평가사 1차 문제/제27회 감정평가사 2교시문제(A형).pdf
%% conversion: kordoc 4.24.51 -> tools/gen1.py (auto); review=proofread-sonnet; flags=none
\begin{question}{65}
(주)감평은 20×1년 초에 해양구조물을 ￦4,000,000(내용연수 5년, 잔존가치 없음, 정액법 상각)에 취득하여 사용하고 있다. 동 해양구조물은 사용기간 종료시점에 원상복구해야 할 의무가 있으며, 종료시점의 원상복구예상금액은 ￦500,000으로 추정되었다. 원가모형을 적용할 경우 (주)감평이 동 해양구조물의 회계처리와 관련하여 20×1년도 포괄손익계산서에 비용으로 처리할 총 금액은? (단, 유효이자율은 연 10\%이며 단일금액 ￦1의 현가계수(5년, 10\%)는 0.6209이다.)
\begin{choices}
  \item ￦800,000
  \item ￦831,046
  \item ￦862,092
  \item ￦893,135
  \item ￦900,000
\end{choices}
\end{question}
\answer{4}
